\setcounter{section}{1}
\section{Preliminaries}
For a natural number $n \ge 1$ we let $[n] = \{1, \ldots, n\}$.  For a real number $x$ we let $\lfloor x \rceil$ denote the closest integer to $x$.

\subsection{Graph basics and notation}
Let $V$ be a finite set and $V^{(2)}$ the set of all subsets of $V$ of cardinality $2$.  We represent a weighted undirected 
graph as a pair $G = (V,w)$ where $w: V^{(2)} \rightarrow \R$ is a non-negative function.  We let $V(G)$ be the vertex set of a graph $G$ and 
$E(G) = \{ e \in V^{(2)}: w(e) > 0\}$ be the set of edges of $G$.  We extend the weight function to sets $S \subseteq V^{(2)}$ by $w(S) = \sum_{e \in S} w(e)$.
We say that $G$ is \emph{simple} if $w: V^{(2)} \rightarrow \{0,1\}$ and in this case also denote $G$ as $G = (V,E)$, where $E$ is the set of edges.  We call the 
ratio of the largest edge weight of $G$ to the smallest the \emph{edge-weight ratio} of $G$.


For a subset $X \subseteq V$ we use the shorthand $\overline{X} = V \setminus X$, and we say $X$ is \emph{non-trivial} if $\emptyset \ne X \subsetneq V$.  
For disjoint sets $X,Y \subseteq V$ we use $E(X,Y)$ for the set of edges with one endpoint in $X$ and one endpoint in $Y$.
For a non-trivial set $X$, let $\Delta_G(X) = \{ \{i,j\} \in E(G) : i \in X, j \in \overline{X}\}$ be the set of edges 
of $G$ with one endpoint in $X$ and one endpoint in $\overline{X}$.  A \emph{cut} of $G$ is a set of the form $\Delta_G(X)$ for some non-trivial 
set $X$.  We call $X$ and $\overline{X}$ the \emph{shores} of the cut $\Delta_G(X)$.  We call a cut of the form $\Delta_G(\{u\})$ a \emph{star} cut, and refer 
to all other cuts as \emph{non-star} cuts.  The \emph{weight} of a cut $S$ is $w(S)$, which in the 
case of a simple graph equals $|S|$.  We let $\lambda(G) = \min_{\emptyset \ne X \subsetneq V} w(\Delta_G(X))$ be the minimum weight of a cut in $G$.  
We call a cut realizing this bound a \emph{minimum cut}.  We call a cut $\Delta_G(X)$ satisfying $w(\Delta_G(X)) \le \alpha \lambda(G)$ an $\alpha$-near minimum cut. 
In the case where $G$ is simple we call $\lambda(G)$ the \emph{edge connectivity} of $G$.  
We will only use the term edge connectivity in the context of unweighted graphs.


